\documentclass[12pt,a4paper,reqno]{amsart}

% language
\usepackage[greek,english]{babel}
\usepackage[utf8]{inputenc}
\usepackage{alphabeta}

% change default names to greek
\addto\captionsenglish{
    \renewcommand{\contentsname}{Περιεχόμενα}
    \renewcommand{\refname}{Βιβλιογραφία}
    \renewcommand{\datename}{Ημερομηνία:}
    \renewcommand{\urladdrname}{Ιστοσελίδα}
}

% math 
\usepackage{amsmath,amsthm,amssymb,amscd}

% font
\usepackage[cal=euler]{mathalfa}
\usepackage{libertinus-type1}
% \usepackage{txfonts} % for upright greek letters
\usepackage{bm} % for bold symbols
\usepackage{bbm} % for the simply-looking bb symbols

% miscellaneous 
\usepackage{changepage} %for indenting environments
\usepackage{csquotes} % example: \textcquote{}
\usepackage{blkarray}
\setcounter{MaxMatrixCols}{20} % default for pmatrix is 10!!
\usepackage{ytableau}
\usepackage{array} %needed to increase the vertical length in a tabular

% drawing
\usepackage{tikz,tikz-cd}
\usetikzlibrary{shapes.misc, patterns, matrix, calc, intersections,positioning,arrows.meta}
\usepackage{graphics,graphicx}
\usepackage{float} % provides enhanced control and customization options for floating objects such as figures and tables

% colors
\usepackage{xcolor}
\definecolor{darkcandyapplered}{rgb}{0.64, 0.0, 0.0}
\definecolor{midnightblue}{rgb}{0.1, 0.1, 0.44}
\definecolor{mylightblue}{HTML}{336699}
\definecolor{burntorange}{rgb}{0.8, 0.33, 0.0}
\definecolor{iceberg}{rgb}{0.44, 0.65, 0.82}
\definecolor{applegreen}{rgb}{0.55, 0.71, 0.0}
\definecolor{canaryyellow}{rgb}{1.0, 0.94, 0.0}
\definecolor{lightgray}{rgb}{0.83, 0.83, 0.83}

% hrefs
\usepackage{hyperref}
\usepackage[noabbrev,capitalize]{cleveref}
\hypersetup{
    pdftoolbar=true,        
    pdfmenubar=true,        
    pdffitwindow=false,     
    pdfstartview={FitH},  % fits the width of the page to the window
    pdftitle={},
    pdfauthor={},
    pdfsubject={},
    pdfkeywords={},
    pdfnewwindow=true,  % links in new window
    colorlinks=true,  % false: boxed links; true: colored links
    linkcolor=darkcandyapplered,   % color of internal links
    citecolor=midnightblue,  % color of links to bibliography
    urlcolor=cyan,  % color of external links
    linktocpage=true  % changes the links from the section body to the page number
    }

% geometry
\textwidth=16cm 
\textheight=21cm 
\hoffset=-55pt 
\footskip=25pt

% thm envs (you might need to change the path)
\theoremstyle{definition}
\newtheorem{theorem}{Θεώρημα}
\newtheorem{proposition}[theorem]{Πρόταση}
\newtheorem{lemma}[theorem]{Λήμμα}
\newtheorem{corollary}[theorem]{Πόρισμα}
\newtheorem{conjecture}[theorem]{Εικασία}
\newtheorem{problem}[theorem]{Πρόβλημα}
\newtheorem*{claim}{Ισχυρισμός}
\newtheorem{observation}[theorem]{Παρατήρηση}
\newtheorem{definition}[theorem]{Ορισμός}
\newtheorem{question}[theorem]{Ερώτημα}
\newtheorem*{questions}{Ερωτήματα}
\newtheorem*{que}{Ερώτημα}
\newtheorem{example}[theorem]{Παράδειγμα}
\newtheorem{exercise}{Άσκηση}

\theoremstyle{remark}
\newtheorem*{remark}{Παρατήρηση}

% fixes the correct numbering of environments
\numberwithin{theorem}{section}
\numberwithin{exercise}{section}
\numberwithin{equation}{section}

% math ops 
% fields
\newcommand{\NN}{\mathbbmss{N}} 
\newcommand{\ZZ}{\mathbbmss{Z}} 
\newcommand{\QQ}{\mathbbmss{Q}} 
\newcommand{\RR}{\mathbbmss{R}} 
\newcommand{\CC}{\mathbbmss{C}} 
\newcommand{\KK}{\mathbbmss{K}} 
\newcommand{\FF}{\mathbbmss{F}} 
\newcommand{\HH}{\mathbbmss{H}} 

% symmetric group
\newcommand{\fS}{\mathfrak{S}}  

% calligraphic 
\newcommand{\aA}{\mathcal{A}} 
\newcommand{\bB}{\mathcal{B}}
\newcommand{\cC}{\mathcal{C}}
\newcommand{\dD}{\mathcal{D}}
\newcommand{\eE}{\mathcal{E}}
\newcommand{\fF}{\mathcal{F}}
\newcommand{\hH}{\mathcal{H}}
\newcommand{\iI}{\mathcal{I}}
\newcommand{\lL}{\mathcal{L}}
\newcommand{\oO}{\mathcal{O}}
\newcommand{\pP}{\mathcal{P}}
\newcommand{\sS}{\mathcal{S}}
\newcommand{\mM}{\mathcal{M}}
\newcommand{\uU}{\mathcal{U}}

% bold
\newcommand{\bfa}{\mathbf{a}}
\newcommand{\bfe}{\mathbf{e}}
\newcommand{\bfF}{\pmb{F}}
\newcommand{\bfR}{\pmb{R}}
\newcommand{\bfv}{\mathbf{v}}
%\newcommand{\bfx}{\bm{x}}
%\newcommand{\bfx}{\mathbf{x}} 
\newcommand{\bfx}{\pmb{x}}
\newcommand{\bfX}{\pmb{X}}
\newcommand{\bfy}{\pmb{y}}
\newcommand{\bfz}{\pmb{z}}

% roman
\newcommand{\rmA}{\mathrm{A}}
\newcommand{\rmB}{\mathrm{B}}
\newcommand{\rmC}{\mathrm{C}}
\newcommand{\rmD}{\mathrm{D}} 
\newcommand{\rmH}{\mathrm{H}} 
\newcommand{\rmh}{\mathrm{h}} 
\newcommand{\rmI}{\mathrm{I}} 
\newcommand{\rmK}{\mathrm{K}}
\newcommand{\rmM}{\mathrm{M}}
\newcommand{\rmP}{\mathrm{P}}  
\newcommand{\rmp}{\mathrm{p}}  
\newcommand{\rmQ}{\mathrm{Q}}  
\newcommand{\rmR}{\mathrm{R}}
\newcommand{\rmS}{\mathrm{S}}
\newcommand{\rmT}{\mathrm{T}}
\newcommand{\rmU}{\mathrm{U}}
\newcommand{\rmV}{\mathrm{V}}
\newcommand{\rmY}{\mathrm{Y}}
\newcommand{\rmZ}{\mathrm{Z}}
\newcommand{\rmz}{\mathrm{z}}

% arrows and symbols 
\renewcommand{\to}{\rightarrow}
\newcommand{\toto}{\longrightarrow}
\newcommand{\mapstoto}{\longmapsto}
\newcommand{\then}{\Rightarrow}
\newcommand{\IFF}{\Leftrightarrow}
\newcommand{\tl}{\tilde}
\newcommand{\wtl}{\widetilde}
\newcommand{\ol}{\overline}
\newcommand{\ul}{\underline}
\newcommand{\oldemptyset}{\emptyset}
\renewcommand{\emptyset}{\varnothing}
\DeclareMathSymbol{\Arg}{\mathbin}{AMSa}{"39} % for arguments 
\newcommand{\onto}{\ensuremath{\twoheadrightarrow}}
\newcommand{\tle}{\trianglelefteq}
\newcommand{\tge}{\trianglerighteq}

% custom ops
\newcommand{\rmKer}{\mathrm{Ker}}
\newcommand{\rmIm}{\mathrm{Im}}
\newcommand{\lcm}{\mathrm{lcm}}
\newcommand{\GL}{\mathrm{GL}}
\newcommand{\ev}{\mathrm{ev}}
\newcommand{\End}{\mathrm{End}}
\newcommand{\rmid}{\mathrm{id}}
\newcommand{\rmspan}{\mathrm{span}}
\newcommand{\Hom}{\mathrm{Hom}}
\newcommand{\Ann}{\mathrm{Ann}}

% absolute value symbol
\usepackage{mathtools} 
\DeclarePairedDelimiter\abs{\lvert}{\rvert}%
\DeclarePairedDelimiter\norm{\lVert}{\rVert}%
\makeatletter
\let\oldabs\abs
\def\abs{\@ifstar{\oldabs}{\oldabs*}}

% tensor symbol
\newcommand{\tensor}[1]{%
  \mathbin{\mathop{\otimes}\limits_{#1}}%
}

% permutation cycle notation
\ExplSyntaxOn
\NewDocumentCommand{\cycle}{ O{\;} m }
 {
  (
  \alec_cycle:nn { #1 } { #2 }
  )
 }

\seq_new:N \l_alec_cycle_seq
\cs_new_protected:Npn \alec_cycle:nn #1 #2
 {
  \seq_set_split:Nnn \l_alec_cycle_seq { , } { #2 }
  \seq_use:Nn \l_alec_cycle_seq { #1 }
 }
\ExplSyntaxOff

% setminus symbol
\newcommand{\mysetminusD}{\hbox{\tikz{\draw[line width=0.6pt,line cap=round] (3pt,0) -- (0,6pt);}}}
\newcommand{\mysetminusT}{\mysetminusD}
\newcommand{\mysetminusS}{\hbox{\tikz{\draw[line width=0.45pt,line cap=round] (2pt,0) -- (0,4pt);}}}
\newcommand{\mysetminusSS}{\hbox{\tikz{\draw[line width=0.4pt,line cap=round] (1.5pt,0) -- (0,3pt);}}}
\newcommand{\sm}{\mathbin{\mathchoice{\mysetminusD}{\mysetminusT}{\mysetminusS}{\mysetminusSS}}}

% miscellaneous commands
\newcommand{\defn}[1]{{\color{mylightblue}{#1}}}
\newcommand{\toDo}{{\bf\color{red} TODO}}
\newcommand{\toCite}{{\bf\color{green} CITE}}
\newcommand*{\vertbar}{\rule[-1ex]{0.5pt}{2.5ex}} % for matrices with column vectors
\newcommand*{\horzbar}{\rule[.5ex]{2.5ex}{0.5pt}} % for matrices with row vectors
\newcommand{\myblue}[1]{{\color{iceberg}{#1}}}
\newcommand{\myorange}[1]{{\color{burntorange}{#1}}}
\newcommand{\mygreen}[1]{{\color{applegreen}{#1}}}
\newcommand{\myred}[1]{{\color{darkcandyapplered}{#1}}}

% ferrer's diagram
\newcommand{\fdiagram}[1]{
    \begin{tikzpicture}[scale=.7]
        \fill foreach \Z [count=\Y] in {#1}
        {foreach \X in {1,...,\Z} 
        {(\X,-\Y) circle[radius=3pt]}};
    \end{tikzpicture}
}

%
\usepackage{pgfplots}
\pgfplotsset{compat=1.18}

%
\makeatletter
\def\mathcolor#1#{\@mathcolor{#1}}
\def\@mathcolor#1#2#3{%
  \protect\leavevmode
  \begingroup
    \color#1{#2}#3%
  \endgroup
}
\makeatother

% 
\newenvironment{nouppercase}{%
  \let\uppercase\relax%
  \renewcommand{\uppercasenonmath}[1]{}}{}

% titlepage
\title{226: Θεωρία Δακτυλίων και Modules}
\author[Β.~Δ. Μουστακας]{Βασίλης Διονύσης Μουστάκας \\ Πανεπιστήμιο Κρήτης}
\date{10 Μαρτίου 2026}
% \urladdr{\href{https://sites.google.com/view/vasmous}{https://sites.google.com/view/vasmous}}

\begin{document}

\begingroup
\def\uppercasenonmath#1{} % this disables uppercase title
\let\MakeUppercase\relax % this disables uppercase authors
\maketitle
\endgroup

\setcounter{section}{4}
\setcounter{theorem}{5}
% \setcounter{equation}{}
\begin{center}
    \textbf{4. Πρότυπα: Εισαγωγή
} (Συνέχεια)
\end{center}

Υπενθυμίζουμε πως σε ότι ακολουθεί, $R$ είναι ένας (όχι κατά ανάγκη μεταθετικός) δακτύλιος.

\begin{definition}
    \label{def:module_homomorphism}
    Έστω $M, N$ δύο $R$-πρότυπα. 
    \begin{itemize}
        \item Μία απεικόνιση $\phi : M \to N$ ονομάζεται \defn{ομομορφισμός $R$-προτύπων} (ή \defn{$R$-ομομορφι\-σμός} ή \defn{$R$-γραμμική}) αν διατηρεί 
        \begin{itemize}
            \item το μηδέν, δηλαδή $\phi(0_M) = 0_N$
            \item την πρόσθεση, δηλαδή $\phi(a + b) = \phi(a) + \phi(b)$
            \item την δράση του $R$, δηλαδή $\phi(rm) = r\phi(m)$,
        \end{itemize}
        για κάθε $a, b, m \in M$ και $r \in R$.
        \item Μια $R$-γραμμική απεικόνιση $M \to N$ η οποία είναι 1-1 και επί ονομάζεται \defn{$R$-ισομο\-ρφισμός} και στην περίπτωση αυτή γράφουμε $M \cong N$.
        \item Γράφουμε $\Hom_R(M,N)$ για το σύνολο όλων των $R$-γραμμικών απεικονίσεων $M \to N$ και θέτουμε 
        \[
        \End_R(M) \coloneqq \Hom_R(M,M).
        \]
    \end{itemize}
\end{definition}

Αν ο $R$ είναι σώμα (αντ. $\ZZ$), τότε οι $R$-γραμμικές απεικονίσεις, είναι απλώς οι γραμμικές απεικονίσεις (αντ. ομομορφισμοί αβελιανών ομάδων). 

\begin{example}
    \label{ex:module_homomorphism}
    \leavevmode
    \begin{itemize}
        \item[(α)] Για κάθε θετικό ακέραιο $n$, η απεικόνιση 
        \begin{align*}
            \phi_n : \ZZ &\to \ZZ \\ 
            a &\mapsto na 
        \end{align*}
        είναι $\ZZ$-γραμμική. Αυτή είναι ομομορφισμός δακτυλίων αν και μόνο αν $n=1$.
        \item[(β)] Γενικότερα, αν $R$ είναι μεταθετικός δακτύλιος και $M$ είναι $R$-πρότυπο, τότε για κάθε $r \in R$ η απεικόνιση 
        \begin{align*}
            M &\to M \\ 
            a &\mapsto ra 
        \end{align*}
        είναι $R$-γραμμική (γιατί;).
        \item[(γ)] Για κάθε $1 \le i \le n$, η προβολή στην $i$-οστή συντεταγμένη του $R^n$ 
        \begin{align*}
            \pi_i : R^n &\to R \\ 
            (r_1, r_2, \dots, r_n) &\mapsto r_i
        \end{align*}
        είναι $R$-επιμορφισμός
        \item[(δ)] Αν ο $R$ είναι μεταθετικός δακτύλιος, τότε η απεικόνιση 
        \begin{align*}
             \rmM_{n,m}(R) &\to \rmM_{m,n}(R) \\ 
            A &\mapsto A^\top
        \end{align*}
        όπου $A^\top$ είναι ο \emph{ανάστροφος} του $A$ είναι $R$-ισομορφισμός (γιατί;).
        \item[(ε)] Η απεικόνιση 
        \begin{align*}
            \{(r_1, r_2, \dots, r_n) \in R^n : r_1 = r_2 = \cdots = r_n\} &\to R \\ 
            (r_1, r_2, \dots, r_n) &\mapsto r_1
        \end{align*}
        είναι $R$-ισομορφισμός.
    \end{itemize}
\end{example}

\begin{proposition}
    \label{prop:Hom}
    Έστω $M, N$ δύο $R$-πρότυπα. 
    \begin{itemize}
        \item Το $\Hom_R(M,N)$ εφοδιασμένο με την πρόσθεση που ορίζεται θέτοντας
        \[
        (\phi + \psi)(m) \coloneqq \phi(m) + \psi(m) 
        \]
        για κάθε $m \in M$ και για κάθε $\phi, \psi \in \Hom_R(M,N)$ και μηδέν την μηδενική απεικόνιση είναι αβελιανή ομάδα.
        \item Επιπλέον, αν ο $R$ είναι μεταθετικός δακτύλιος, τότε το $\Hom_R(M,N)$ γίνεται $R$-πρότυ\-πο με την δράση του $R$ που ορίζεται θέτοντας
        \[
        \left(r\cdot\phi\right)(m) \coloneqq r\cdot\phi(m)
        \]
        για κάθε $r \in R$ και $\phi \in \Hom_R(M,N)$.
        \item Το $\End_R(M)$ γίνεται δακτύλιος με πολλαπλασιασμό την σύνθεση απεικονίσεων και μονάδα την ταυτοτική απεικόνιση, ο οποίος ονομάζεται \defn{δακτύλιος ενδομορφισμών} του $M$.
    \end{itemize}
\end{proposition}

\begin{proof}[Απόδειξη]
    Η απόδειξη αφήνεται ως άσκηση. Παρατηρούμε μόνο ότι, στον δεύτερο ισχυρισμό, για είναι η δράση του $R$ στο $\Hom_R(M,N)$ καλώς ορισμένη, δηλαδή να ισχύει ότι $r\cdot \phi \in \Hom_R(M,N)$ πρέπει
    \[
    \left(r\cdot\phi\right)(sm) = s\left(r\cdot\phi\right)(m)
    \]
    το οποίο είναι ισοδύναμο με 
    \[
    rs\phi(m) = sr\phi(m)
    \]
    το οποίο με τη σειρά του είναι ισοδύναμο με 
    \[
    rs = sr
    \]
    για κάθε $r, s \in R$, απ' όπου προκύπτει η υπόθεση ο $R$ να είναι μεταθετικός δακτύλιος.
\end{proof}

Η $R$-γραμμική απεικόνιση του Παραδείγματος \ref{ex:module_homomorphism} (β) επάγει έναν ομομορφισμό δακτυλίων 
\[
R \to \End_R(M)
\]
με πυρήνα\footnote{Το ιδεώδες αυτό ονομάζεται \emph{μηδενιστής} (annihilator) του $M$ και συνήθως συμβολίζεται με $\Ann_R(M)$.}
\[
\{r \in R : rm = 0_M, \ \text{για κάθε $m \in M$}\}.
\]

\begin{definition}
    \label{def:kernel_image}
    Έστω $M, N$ δύο $R$-πρότυπα και $\phi \in \Hom_R(M,N)$. Τα σύνολα 
    \begin{align*}
        \ker(\phi) &\coloneqq \{a \in M : \phi(a) = 0_M\} \\
        \rmIm(\phi) &\coloneqq \{b \in M : b = \phi(a), \ \text{για κάποιο $a \in M$}\} 
    \end{align*}
    ονομάζονται \defn{πυρήνας} και \defn{εικόνα} της $\phi$, αντίστοιχα.
\end{definition}

Ο πυρήνας και η εικόνα ενός $R$-ομομορφισμού $\phi : M \to N$ είναι υποπρότυπα των $M$ και $N$, αντίστοιχα. Επιπλέον, η $\phi$ είναι $R$-ισομορφισμός αν και μόνο αν 
\[
\ker(\phi) = \{0_m\} \quad \text{και} \quad \rmIm(\phi) = N.
\]

Για παράδειγμα, για τον πολλαπλασιασμό με κάποιο θετικό ακέραιο $n$, τον οποίο είδαμε στο Παράδειγμα \ref{ex:module_homomorphism} (α), υπολογίζουμε
\[
\ker(\phi_n) = \{0\} \quad \text{και} \quad \rmIm(\phi_n) = n\ZZ.
\]
Μπορούμε να συνοψίσουμε την πληροφορία που μας δίνει η $\phi_n$ στο εξής διάγραμμα $\ZZ$-προτύπων\footnote{Με $\hookrightarrow$ και $\twoheadrightarrow$ συνήθως συμβολίζουμε τους μονομορφισμούς και τους επιμορφισμούς, αντίστοιχα.}
\begin{equation} 
    \label{eq:ses_Z}
    \begin{tikzcd}
    0 \arrow[r] & \mathbb{Z} \arrow[r, hook, "\phi_n"] & \mathbb{Z} \arrow[r, two heads,"\pi"] & \mathbb{Z}_n \arrow[r] & 0
    \end{tikzcd}
\end{equation}
όπου η $\pi : \ZZ \to \ZZ_n$ είναι ο \textquote{φυσικός} επιμορφισμός 
\[
\pi(a) = a + n\ZZ
\]
και με $0$ συμβολίζουμε το τετριμμένο $\ZZ$-πρότυπο\footnote{Παρατηρήστε ότι δε προσδιορίσαμε τους ομομορφισμούς από και προς το τετριμμένο πρότυπο, διότι υπάρχει μόνο μία επιλογή στην ουσία (ποιά;)}.
Παρατηρούμε ότι 
\begin{itemize}
    \item η $\phi_n$ είναι $\ZZ$-μονομορφισμός,
    \item η $\pi$ είναι $\ZZ$-επιμορφισμός, και 
    \item $\ker(\pi) = \rmIm(\phi_n) = n\ZZ$.
\end{itemize}

\begin{definition}
    \label{def:ses}
    Μια ακολουθία $R$-προτύπων 
    \[
    \begin{tikzcd}
        X \arrow[r, "\alpha"] & Y \arrow[r, "\beta"] & Z 
    \end{tikzcd}
    \]
    ονομάζεται \defn{ακριβής} (exact) στο $Y$ αν $\ker(\beta) = \rmIm(\alpha)$. Γενικότερα, μια ακολουθία $R$-προτύπων 
    \[
    \begin{tikzcd}
        \cdots \arrow[r] & X_{n-1} \arrow[r, "\phi_{n-1}"] & X_n \arrow[r, "\phi_n"] & X_{n+1} \arrow[r] & \cdots 
    \end{tikzcd}
    \]
    ονομάζεται \defn{ακριβής} αν είναι ακριβής σε κάθε $X_n$.
\end{definition}

Παρατηρούμε ότι 
\begin{itemize}
    \item Η ακολουθία 
    \[
    \begin{tikzcd}
        0 \arrow[r] & X \arrow[r, "\alpha"] & Y 
    \end{tikzcd}
    \]
    είναι ακριβής στο $X$ αν και μόνο αν η $\alpha$ είναι $R$-μονομορφισμός (γιατί;) και 
    \item Η ακολουθία 
    \[
    \begin{tikzcd}
         Y \arrow[r, "\beta"] & Z \arrow[r] & 0
    \end{tikzcd}
    \]
    είναι ακριβής στο $Z$ αν και μόνο αν η $\beta$ είναι $R$-επιμορφισμός (γιατί;).
\end{itemize}
Συνεπώς, η ακολουθία $\ZZ$-προτύπων \eqref{eq:ses_Z} είναι ακριβής και έχει την μορφή 
\[
\begin{tikzcd}
        0 \arrow[r] & X \arrow[r, "\alpha"] & Y \arrow[r, "\beta"] & Z \arrow[r] & 0.
\end{tikzcd}
\]
Οι ακριβείς ακολουθίες αυτής της μορφής ονομάζονται \defn{βραχείες ακριβείς ακολουθίες} (short exact sequences).

\begin{proposition}
    \label{prop:quotient_module}
    Έστω $M$ ένα $R$-πρότυπο και $N$ ένα υποπρότυπο του $M$. Η αβελιανή ομάδα πηλίκο $M/N$ εφοδιασμένη με τη δράση του $R$ που ορίζεται θέτοντας 
    \[
    r \cdot{m + N} \coloneqq rm + N
    \]
    για κάθε $r \in R$ και $m \in M$ έχει δομή $R$-προτύπου και ονομάζεται \defn{πρότυπο πηλίκο}. Επιπλέον, η απεικόνιση 
    \begin{align*}
            \pi : M &\to M/N \\ 
                  m &\mapsto m + N 
    \end{align*}
    είναι $R$-επιμορφισμός με πυρήνα $N$.
\end{proposition}

Η απόδειξη της Πρότασης \ref{prop:quotient_module} αφήνεται ως άσκηση. Αν ο $R$ είναι σώμα, τότε τα πρότυπα πηλίκο είναι απλώς οι (γραμμικοί) χώροι πηλίκο. Στην περίπτωση όπου $R=\ZZ$, τα πρότυπα πηλίκο είναι απλώς οι (αβελιανές) ομάδες πηλίκο.

\begin{theorem}{\rm(Θεωρήματα Ισομορφισμών)}
    \label{thm:isomorphism_theorems}
    \leavevmode
    \begin{itemize}
        \item[(1)] Αν $\phi : M \to N$ είναι $R$-ομομορφισμός, τότε $M/\ker(\phi) \cong \rmIm(\phi)$.
        \item[(2)] Αν $A, B$ είναι υποπρότυπα ενός $R$-πρτύπου $M$, τότε $A/(A\cap{B}) \cong (A+B)/B$, όπου $A + B \coloneqq \{a + b : a \in A, b \in B\}$.
        \item[(3)] Αν $N$ και $L$ είναι δύο υποπρότυπα ενός $R$-προτύπου $M$ τέτοια ώστε $L \subseteq N$, τότε $\frac{M/L}{N/L} \cong M/N$.
        \item[(4)] Έστω $M$ ένα $R$-πρότυπο. Υπάρχει μια αμφιμονοσήμαντη αντιστοιχία μεταξύ των υποπροτύπων του $M/N$ και των υποπροτύπων του $N$ τα οποία περιέχουν το $N$. 
    \end{itemize}
\end{theorem}

Η απόδειξη του Θεωρήματος \ref{thm:isomorphism_theorems} είναι όμοια με αυτή του Θεωρήματος 1.9.
\begin{remark}
    Κάθε $R$-ομομοριφμός $\phi : M \to N$ επάγει μία βραχεία ακριβή ακολουθία $R$-προτύπων
    \[
    \begin{tikzcd}
        0 \arrow[r] & \ker(\phi) \arrow[r, hook] & M \arrow[r, two heads] & \rmIm(\phi) \arrow[r] & 0
    \end{tikzcd}
    \]
    (γιατί;). Το πρώτο θεώρημα ισομορφισμού μας πληροφορεί ότι 
    $\rmIm(\phi) \cong M/\ker(\phi)$.

    Αντιστρόφως, δοθείσης μίας βραχείας ακριβής ακολουθίας $R$-προτύπων
    \[
    \begin{tikzcd}
        0 \arrow[r] & X \arrow[r, "\alpha"] & Y \arrow[r, "\beta"] & Z \arrow[r] & 0
    \end{tikzcd}
    \]
    έχουμε τον εξής $R$-ισομορφισμό
    \[
    Z \cong Y/\alpha(X).
    \]
    Πράγματι, 
    \[
    Y/\alpha(X) \cong Y/\ker(\beta) \cong \beta(Y) \cong Z,
    \]
    όπου ο πρώτος ισομορφισμός έπεται από την ακρίβεια στο $X$, ο δεύτερος από το πρώτο θεώρημα ισομορφισμού και ο τρίτος από την ακρίβεια στο $Z$.

    Με άλλα λόγια, μπορούμε να σκεφτόμαστε κάθε βραχεία ακριβή ακολουθία ως ένα (πρώ\-το) θεώρημα ισομορφιμού. Για παράδειγμα, η βραχεία ακριβής ακολουθία $\ZZ$-προτύπων \eqref{eq:ses_Z} \textquote{γνωρίζει} την πληροφορία 
    \[
    \ZZ_n \cong \ZZ/\phi_n(\ZZ) = \ZZ/n\ZZ.
    \]
\end{remark}

% \begin{thebibliography}{99}
%     \bibitem{DF04}
%     D. S. Dummit και R. M. Foote, 
%     \textit{Abstract algebra},
%     third edition, Wiley, 2004. 
% \end{thebibliography}
\end{document}